\scanpage{70}
Pour décrire l'action d'un $u\in\unsure{\mathbf{\Gamma}|\Gamma}$ (resp $u\in\Gamma$) sur $\widehat\Pi_1(X)$ (resp. $\widehat\Pi_1(X_{\bar K})$), on prendra comme ``générateurs'' du groupoïde $\widehat\Pi_1$ des chemins provenant de $\Pi_1$. On aura, pour ${l}:x\to y$ dans $\Pi_1$, ${}^u{l}:{}^ux\to{}^uy$.

Si \deleted{on suppose} par la suite, pour simplifier, on se borne à des points $x,y\ldots$ rationnels sur $\unsure{X|K}$ -- i.e. qu'on travaille dans le sous-groupoïde $\widehat\Pi_1(X_K)$ formé des objets qui sont invariants sous $\unsure{\mathbf{\Gamma}|\Gamma}$ -- alors on aura
\[
(*)\qquad{}^u{l}:{}^ux=x\longrightarrow{}^uy=y,
\]
et on peut écrire, \deleted{\unsure{f}}
\[
(8)\qquad\begin{cases}
{}^u{l}=g_{l}(u)\circ{l}={l}\circ f_{l}(u)\\
g_{l}(u)\in\widehat\pi_1(X,y)\\
f_{l}(u)\in\widehat\pi_1(X,x)
\end{cases}
\]
\deleted{$g_{l}(u)={}^u{l}\circ{l}^{-1}$}
\figure{\begin{tikzpicture}[x=2mm,y=2mm]
\node at(0,0){,};
\fill(.5,.8)circle(.07);\fill(.75,.55)circle(.07);\fill(1.25,.4)circle(.08);\fill(1.75,.08)circle(.045);
\fill[gray](1.16,-.08)circle(.04);\fill[gray](1.38,-.1)circle(.035);
\end{tikzpicture}}
\margin{On écrira aussi parfois $g_u({l})$ au lieu de $g_{l}(u)$, $f_u({l})$ pour $f_{l}(u)$.}
où $f_{l}(u),g_{l}(u),{}^u{l}$ se déterminent mutuellement, \deleted{\unsure{fonctoriquement}} par les formules (8) et
\[
(9)\qquad\begin{cases}
f_{l}(u)={l}^{-1}\cdot{}^u{l},\quad g_{l}(u)={}^u{l}\cdot{l}^{-1}\\
g_{l}(u)={l}(f_{l}(u)),\quad f_{l}(u)={l}^{-1}(g_{l}(u))
\end{cases}
\]
On décrira l'action d'un $u\in\Gamma$ \inserted{sur un sous-groupoïde plein de $\Pi_1(X)$} à l'aide des éléments $f_{l}(u)$, ou $g_{l}(u)$, pour un syst. de générateurs convenable des \inserted{groupoïdes pleins} correspondants de $\Pi_1(X)$. Dans le cas d'un
\margin{69}
